\documentclass[12pt,a4paper]{article}

\usepackage[utf8]{inputenc}
\usepackage[english]{babel}
\usepackage{amsmath,amssymb}
\usepackage{graphicx}
\usepackage{booktabs}
\usepackage{hyperref}
\usepackage{natbib}
\usepackage{float}
\usepackage{multirow}
\usepackage{caption}
\usepackage{subcaption}
\usepackage{colortbl}
\usepackage{xcolor}
\usepackage[margin=2.5cm]{geometry}
\usepackage{microtype}

\graphicspath{{../../slides_export/}}

\definecolor{bestrow}{RGB}{200,230,200}
\definecolor{greyrow}{RGB}{242,242,242}

\title{Biotic Interaction Detection from Scientific Literature:\\
Knowledge Distillation and Multi-task Learning\\
with Named Entity Recognition}

\author{
Esteban Gaillac\textsuperscript{1,*} \\[4pt]
\textsuperscript{1}Haute \'Ecole de Gestion, HES-SO Geneva \\
\textsuperscript{*}Correspondence: esteban.gaillac@hesge.ch
}

\date{June 2026}

\begin{document}

\maketitle

% ─────────────────────────────────────────────────────────────────────────────
\begin{abstract}
Detecting biotic interactions from scientific text is a prerequisite for
automated knowledge graph construction from biodiversity literature.
We present a multi-task transformer classifier that achieves F1\,=\,0.868
(P\,=\,0.793, R\,=\,0.958) on a 100-sentence ecological test set,
outperforming both a fine-tuned BiomedBERT baseline (F1\,=\,0.788) and
a BiomedBERT\,$\times$\,FLAN-T5-base ensemble (F1\,=\,0.865) at 3.3$\times$
lower inference cost.

The model is built on three components. First, a two-stage teacher chain
in which a local 122-billion-parameter model (Qwen3.5-122B) annotates
44,178 harvested sentences with binary labels, and a trained ensemble
subsequently generates calibrated soft probability labels for knowledge
distillation. Second, a multi-task BiomedBERT model that jointly classifies
sentences and annotates tokens with biological entity roles (HOST, PATHOGEN,
SPECIES, INT) using automatically derived BIO supervision, with a two-phase
training strategy that primes the shared encoder via NER before introducing
the classification objective. Third, we show that soft labels are the single
largest performance lever in the pipeline: switching from hard to soft
cross-entropy on identical architecture and data yields a +0.052 F1 gain.

A distribution analysis reveals that 77\% of test-set positives belong to
interaction types underrepresented in the PubMed-sourced training corpus,
explaining the current performance ceiling at F1\,$\approx$\,0.87 and
providing a clear target for future data collection. The classifier is
deployed upstream of BiotXplorer
(\url{https://biotxplorer.sibils.org/}), the Swiss Institute of Bioinformatics
semantic search platform for biotic interactions, and is designed as the
first component of a downstream knowledge graph construction pipeline.
\end{abstract}

\textbf{Keywords:} biotic interactions, knowledge distillation, multi-task
learning, named entity recognition, BiomedBERT, knowledge graph,
ecological text mining, GloBI

% ─────────────────────────────────────────────────────────────────────────────
\section{Introduction}
\label{sec:intro}

Species interactions shape ecosystem structure, drive disease dynamics,
and determine the response of ecological communities to perturbation.
The Global Biotic Interactions database (GloBI; \citealt{poelen2014global})
aggregates over 13 million structured interaction records from more than 700
data sources, yet the vast majority of interaction knowledge remains encoded
in unstructured scientific literature distributed across millions of
publications. Europe PMC alone indexes more than 40 million articles.
Systematic curation of this corpus is infeasible at human scale.

Sentence-level classification of biotic interactions is the entry point
for automated literature mining. Given a sentence, the classifier determines
whether it describes a biotic interaction between two named organisms:
predation, parasitism, pathogen infection, pollination, herbivory, mutualism,
or symbiosis. This binary decision feeds two downstream applications.

\paragraph{Upstream: BiotXplorer.}
BiotXplorer \citep{sibils2023biotxplorer} is a semantic search engine for
biotic interactions in biomedical literature, maintained by the Swiss Institute
of Bioinformatics. The classifier acts as an upstream quality filter, screening
millions of candidate sentences to retain only those describing real interactions
before they are indexed. A high-recall, precision-conscious filter of this kind
reduces index noise while preserving comprehensive coverage across interaction
types.

\paragraph{Downstream: knowledge graph construction.}
Beyond search, structured interaction triples extracted from accepted sentences
can populate biodiversity knowledge graphs \citep{chen2023biodiversity}
supporting species-level reasoning about disease spread, invasion risk,
food web dynamics, and conservation planning \citep{poisot2021global}.
The classifier is the first component in such a pipeline; downstream relation
extraction maps accepted sentences to typed triples of the form
(subject entity, relation, object entity, confidence, provenance).

Previous NLP work in the biomedical domain has focused primarily on
protein-protein interactions \citep{bunescu2005comparative} and drug-gene
relationships \citep{lim2018drug}, with ecological interaction types
(predation, parasitism, pollination, herbivory) receiving comparatively less
attention. The ecological literature presents additional challenges: it spans
thousands of journals, uses varied terminology, and describes field
observations that differ structurally from the experimental prose dominant
in PubMed \citep{badenes2020detecting}.

We address this gap with three primary contributions:
\begin{enumerate}
    \item A two-stage teacher chain combining large-LLM binary annotation
    with ensemble-based soft probability distillation, yielding a
    44K-sentence training corpus with calibrated labels.
    \item A multi-task BiomedBERT architecture with automatic NER auxiliary
    supervision, trained in two phases to prime the shared encoder.
    \item A quantitative distribution mismatch analysis identifying the
    interaction types responsible for the current performance ceiling and
    guiding future data collection.
\end{enumerate}

% ─────────────────────────────────────────────────────────────────────────────
\section{Methods}
\label{sec:methods}

\subsection{Task Definition}

Given a sentence $s$ from scientific literature, predict $y \in \{0, 1\}$
where $y=1$ indicates that $s$ describes a biotic interaction between two
or more named organisms. Interaction types in scope: predation, parasitism,
pathogen infection, pollination, herbivory, mutualism, symbiosis, seed
dispersal, competition, and vector-borne transmission.

\subsection{Training Corpus Construction}
\label{sec:data}

\paragraph{Sentence sources.}
Two sentence origins are combined into a single 44,178-sentence corpus.

\emph{GloBI-derived templates (v1--v7, synthetic).}
Structured triples from GloBI (e.g.\ \textit{parasitizes(Anisakis simplex,
Gadus morhua)}) are injected into hand-written sentence templates covering
nine interaction categories. In the final template version (v7), every
positive sentence was individually validated by the Claude API (Anthropic)
to remove formulaic or linguistically implausible candidates, yielding
25,081 clean sentences (7,251 positive; F1\,=\,0.788).

\emph{Real PMC sentences (v10+).}
Articles co-mentioning known GloBI species pairs were queried from Europe PMC
($\approx$31K sentences) and Biodiversity PMC ($\approx$8K sentences, queried
by interaction category via the SIBiLS API). Co-occurrence sentences were
extracted and labeled by a local Qwen3.5-122B model using the prompt:
\textit{``Does this sentence describe a biotic interaction between two named
organisms? Answer yes or no only.''}
Of 44,178 harvested sentences, 4,065 (9.2\%) were accepted as positives.

\paragraph{Soft label generation.}
The accepted 44,178 sentences were re-scored by the trained BiomedBERT
$\times$ FLAN-T5-base ensemble, producing a soft probability
$p_\text{ensemble} \in [0,1]$ per sentence. These soft labels constitute
the distillation teacher signal (Section~\ref{sec:distillation}).
Figure~\ref{fig:version_f1} summarises F1 on the test set across dataset
configurations.

\begin{figure}[H]
    \centering
    \includegraphics[width=0.85\textwidth]{fig_01_version_f1}
    \caption{Test F1 by training data configuration.
    The best result (rightmost) is the 44K Qwen-labeled corpus with
    ensemble soft labels.}
    \label{fig:version_f1}
\end{figure}

\subsection{Model Architecture Evolution}
\label{sec:architectures}

We evaluated five approaches in order of increasing complexity.

\paragraph{Classical baselines.}
SVM and Random Forest on TF-IDF features achieved F1\,$\approx$\,0.62,
establishing the need for contextual representations.

\paragraph{Fine-tuned BiomedBERT.}
We use \texttt{microsoft/BiomedNLP-BiomedBERT-base-uncased-abstract-fulltext}
\citep{gu2021domain}, a BERT-base model (110M parameters, 12 transformer
layers, 768-dimensional hidden states) pre-trained from scratch on PubMed
abstracts and PubMed Central full-text articles. A linear classification
head on the [CLS] token representation is fine-tuned using 5-fold
cross-validation, AdamW optimiser ($\eta = 2 \times 10^{-5}$, weight
decay\,=\,0.01), and threshold optimisation on the validation split.
BiomedBERT trained on v7 templates achieves F1\,=\,0.788.

\paragraph{BiomedBERT\,$\times$\,FLAN-T5-base ensemble.}
FLAN-T5-base \citep{chung2022scaling}, a sequence-to-sequence model trained
with instruction tuning, was fine-tuned as a binary classifier by prompting
for a yes/no answer and scoring the first generated token. The geometric
mean of BiomedBERT and FLAN-T5-base predicted probabilities, with joint
threshold optimisation, achieves F1\,=\,0.865. The gain is attributable to
model complementarity: BiomedBERT handles precise, template-like language
while FLAN-T5-base generalises better to paraphrased ecological sentences.

\paragraph{Knowledge distillation.}
Section~\ref{sec:distillation}.

\paragraph{Multi-task BiomedBERT with NER.}
Section~\ref{sec:multitask}.

Figure~\ref{fig:model_evolution} summarises the F1 progression across architectures.

\begin{figure}[H]
    \centering
    \includegraphics[width=0.85\textwidth]{fig_02_model_evolution}
    \caption{Test F1 across model architecture stages. Annotated arrows indicate
    the primary driver of each improvement.}
    \label{fig:model_evolution}
\end{figure}

\subsection{Knowledge Distillation}
\label{sec:distillation}

Knowledge distillation \citep{hinton2015distilling} transfers calibrated
uncertainty from a teacher to a student by training on soft probability
distributions rather than hard labels \citep{bucilua2006model}.
A teacher score of 0.7 instructs the student that a sentence is
\emph{probably} positive but not certain; the student is penalised for
predicting 0.4, not merely for predicting the wrong class. This
information transfer is especially valuable for borderline sentences.

\paragraph{Training objective.}
The student BiomedBERT is trained using KL divergence against the ensemble
teacher:
\begin{equation}
    \mathcal{L}_\text{KL} = T^2 \cdot \mathrm{KL}\!\left(
        \sigma(\mathbf{z}_s / T) \;\|\; \sigma(\mathbf{z}_t / T)
    \right)
\end{equation}
where $T$ is the temperature, $\mathbf{z}_s$ and $\mathbf{z}_t$ are student
and teacher logits, and $\sigma$ denotes the softmax function.
Rows without soft labels use hard cross-entropy. The combined loss is
$\mathcal{L} = \alpha \cdot \mathcal{L}_\text{KL} + (1{-}\alpha) \cdot \mathcal{L}_\text{CE}$.

Optimal hyperparameters identified by search: $T{=}2$, $\alpha{=}0.5$
(Table~\ref{tab:distillation}). Temperature $T{=}2$ spreads the teacher
distribution sufficiently to transfer calibrated uncertainty without
over-smoothing; $T{=}4$ over-smooths while $T{=}1.5$ loses the calibration
benefit. The student (distilled\_v2) achieves F1\,=\,0.808, a gain of +0.020
over the v7 fine-tuned baseline.

\begin{table}[H]
\centering
\caption{Knowledge distillation hyperparameter search.
All variants use BiomedBERT as student except v4 (DistilBERT) and v5 (SciBERT).}
\label{tab:distillation}
\begin{tabular}{lccc}
\toprule
\textbf{Version} & $T$ & $\alpha$ & \textbf{Test F1} \\
\midrule
v1 & 4 & 0.7 & 0.785 \\
\rowcolor{bestrow}
\textbf{v2 (best)} & \textbf{2} & \textbf{0.5} & \textbf{0.808} \\
v3 & 4 & 0.9 & 0.808 \\
v4 (DistilBERT) & 2 & 0.5 & 0.792 \\
v5 (SciBERT) & 2 & 0.5 & 0.793 \\
v6 ($T=1.5$) & 1.5 & 0.5 & 0.779 \\
\bottomrule
\end{tabular}
\end{table}

\subsection{Multi-task BiomedBERT with NER Auxiliary Task}
\label{sec:multitask}

\paragraph{Architecture.}
The BiomedBERT encoder is extended with two linear output heads that share
all encoder weights:
\begin{itemize}
    \item \textbf{Classification head}: reads the [CLS] token vector
    ($\mathbb{R}^{768} \to \mathbb{R}^2$), produces $P(\text{interaction})$.
    \item \textbf{NER head}: reads every token vector
    ($\mathbb{R}^{768} \to \mathbb{R}^{|\mathcal{L}|}$), produces BIO label logits.
\end{itemize}

\paragraph{NER label scheme.}
We use the \texttt{full\_typed} scheme with $|\mathcal{L}| = 9$ labels:
\texttt{O} (background), \texttt{B/I-HOST} (hosting organism),
\texttt{B/I-PATHOGEN} (infecting or parasitic agent),
\texttt{B/I-SPECIES} (ambiguous role in symmetric interactions),
\texttt{B/I-INT} (interaction verb or phrase).
Labels are generated automatically from species gazetteers
($\approx$4.2M binomials plus 21K common names, matched via Aho-Corasick
automata \citealt{aho1975efficient}) and an interaction lexicon (591 GloBI
terms plus 50 biomedical vocabulary terms). Gold species annotations from
training metadata override gazetteer lookup when available.

The HOST/PATHOGEN distinction encodes the direction of the relation: who
acts and who is acted upon. This directional labelling is motivated by the
downstream KG construction task, where edge direction is semantically
significant.

\paragraph{Two-phase training.}
Training proceeds in two phases designed to avoid interference between
the NER and classification objectives \citep{liu2019multi}:

\begin{description}
    \item[Phase 1 --- NER pre-train ($E_\text{NER}$ epochs).]
    The classification head is frozen. Only the encoder and NER head receive
    gradient updates from token-level cross-entropy loss
    $\mathcal{L}_\text{NER}$.
    This phase primes the shared encoder to locate species names and
    interaction verbs before the classification signal is introduced.

    \item[Phase 2 --- Joint training ($E_\text{joint}$ epochs).]
    All parameters are unfrozen. Combined loss:
    \[
        \mathcal{L} = \alpha \cdot \mathcal{L}_\text{cls} +
                      (1{-}\alpha) \cdot \mathcal{L}_\text{NER}
    \]
    where $\mathcal{L}_\text{cls}$ is soft KL divergence against ensemble
    soft labels when available, and hard CE otherwise.
\end{description}

We find that initialising Phase 1 from the best fine-tuned BiomedBERT
checkpoint degrades performance (F1\,=\,0.862 vs.\ cold-start 0.868):
NER gradient updates disrupt the task-specific representations built during
classification fine-tuning, and the subsequent joint phase does not fully
recover them. Cold-starting from the pretrained BiomedBERT weights is
therefore preferred.

\subsection{Evaluation Protocol}
\label{sec:eval}

All models are evaluated on a 100-sentence test set (48 positive)
drawn from ecological literature, curated with relaxed annotation criteria
(any biotic interaction type accepted). The set has no overlap with
training data. We report precision, recall, and F1-score.

Bootstrap confidence intervals are estimated with 10,000 resamples.
Pairwise significance is assessed using McNemar's test
\citep{mcnemar1947note} with continuity correction. For each model,
the decision threshold $\tau$ is optimised on the validation split
by sweeping $\tau \in [0.05, 0.95]$ and selecting the value
maximising F1.

% ─────────────────────────────────────────────────────────────────────────────
\section{Results}
\label{sec:results}

\subsection{Multi-task Ablation}

Table~\ref{tab:ablation} presents the full ablation over NER schemes,
loss weights ($\alpha$), and NER pre-training epochs. All configurations
use the 44K distillation corpus with soft labels except the \texttt{hardce}
variants (hard cross-entropy).

\begin{table}[H]
\centering
\small
\caption{Multi-task ablation results on the test set (100 sentences, 48 positives).
NER pt = NER pre-train epochs before joint training.
Baseline comparisons are shown below the separator.}
\label{tab:ablation}
\begin{tabular}{llccrrr}
\toprule
\textbf{Config} & \textbf{NER scheme} & $\alpha$ & \textbf{NER pt} &
\textbf{F1} & \textbf{P} & \textbf{R} \\
\midrule
basic\_a05           & basic       & 0.5 & 0 & 0.841 & 0.763 & 0.938 \\
typed\_a05           & typed       & 0.5 & 0 & 0.852 & 0.767 & 0.958 \\
full\_a03            & full        & 0.3 & 0 & 0.854 & 0.854 & 0.854 \\
full\_typed\_a05     & full\_typed & 0.5 & 0 & 0.847 & 0.746 & 0.979 \\
full\_a05\_ner2      & full        & 0.5 & 2 & 0.824 & 0.778 & 0.875 \\
full\_typed\_a03\_ner2 & full\_typed & 0.3 & 2 & 0.830 & 0.759 & 0.917 \\
\rowcolor{bestrow}
\textbf{full\_typed\_a05\_ner2} & full\_typed & 0.5 & 2 & \textbf{0.868} & \textbf{0.793} & \textbf{0.958} \\
full\_typed (5 joint ep) & full\_typed & 0.5 & 2+5 & 0.816 & 0.764 & 0.875 \\
\midrule
\rowcolor{greyrow}
v12\_hardce          & full\_typed & 0.5 & 2 & 0.816 & 0.764 & 0.875 \\
\rowcolor{greyrow}
distilled\_v2        & ---         & --- & --- & 0.808 & 0.750 & 0.875 \\
\rowcolor{greyrow}
Ensemble BERT$\times$T5 & ---      & --- & --- & 0.865 & ---   & ---   \\
\bottomrule
\end{tabular}
\end{table}

Three consistent findings emerge from the ablation:

\paragraph{1. Soft labels are the largest single lever (+0.052 F1).}
The \texttt{v12\_hardce} configuration uses the same architecture and data
with hard cross-entropy labels and achieves F1\,=\,0.816. The soft-KL variant
reaches 0.868: a difference of +0.052 F1, which is larger than any
architectural change, including the multi-task NER head (+0.021).

\paragraph{2. NER pre-training is necessary.}
Adding 2 NER pre-training epochs improves \texttt{full\_typed} from 0.847
to 0.868. Extending joint training beyond 3 epochs causes performance to
degrade (F1\,=\,0.816), confirming that a short joint phase after encoder
priming is optimal.

\paragraph{3. Richer NER schemes perform better.}
Performance increases monotonically from basic (0.841) through typed (0.852)
and full (0.854) to full\_typed (0.868 with pre-training). The HOST/PATHOGEN
distinction provides directional supervision that builds richer encoder
representations than undifferentiated species tagging.

\subsection{Statistical Significance}

Figure~\ref{fig:ci} shows 95\% bootstrap confidence intervals on the test set.

\begin{figure}[H]
    \centering
    \includegraphics[width=0.65\textwidth]{fig_manuscript_ci}
    \caption{Bootstrap 95\% confidence intervals for three models on the
    test set (n=100, 48 positives). The significance annotation refers to
    McNemar's test between multi-task and BiomedBERT v7 ($p=0.016$).}
    \label{fig:ci}
\end{figure}

McNemar's test confirms that the multi-task model significantly outperforms
the plain fine-tuned BiomedBERT ($p = 0.016$).
The margin over distilled\_v2 is not individually significant ($p = 0.149$,
12 prediction disagreements), reflecting the limited statistical power of
a 100-sentence test set. The 95\% confidence intervals span approximately
$\pm$0.06--0.07 F1; a test set of 500--1000 sentences would substantially
narrow these intervals.

\subsection{Threshold Analysis}

Figure~\ref{fig:threshold} shows precision, recall, and F1 as a function
of the decision threshold on the test set.

\begin{figure}[H]
    \centering
    \includegraphics[width=0.78\textwidth]{fig_threshold_curve}
    \caption{Precision, recall, and F1 versus decision threshold.
    The F1-optimal threshold ($\tau=0.13$) and a precision-oriented
    threshold ($\tau=0.50$) are annotated. Shaded regions indicate
    high-recall and high-precision operating zones.}
    \label{fig:threshold}
\end{figure}

The F1-optimal threshold ($\tau = 0.13$) is lower than typical because
the model is calibrated conservatively: genuine interactions score
predominantly in the range 0.3--0.7 rather than near 1.0.
For applications where false positives entering a knowledge base are
costly, operating at $\tau = 0.50$ reduces false positives from 12 to 9
at the cost of 6 additional missed interactions (P\,=\,0.816, R\,=\,0.833).
The threshold curve makes this precision-recall trade-off explicitly visible
and allows operators to choose the operating point appropriate to their
downstream task.

\subsection{Distribution Mismatch}

Figure~\ref{fig:distribution} reveals a structural mismatch between
training and test interaction type distributions.

\begin{figure}[H]
    \centering
    \includegraphics[width=0.92\textwidth]{fig_distribution_mismatch}
    \caption{Distribution of interaction types in training positives (left,
    n=4,065 Qwen-labeled sentences) versus test set positives (right, n=48).}
    \label{fig:distribution}
\end{figure}

The training corpus, harvested from PubMed-indexed articles, is dominated by
predation (42\%), herbivory (22\%), and pollination (22\%) sentences sourced
via targeted PMC queries. Parasitism/Host interactions represent only 10\%
of training positives despite constituting 42\% of test-set positives.
Pathogen/Infection represents 3\% of training but 35\% of test-set positives.
Together, these two types account for 77\% of test-set positives but only
13\% of training data.

This mismatch arises from a structural property of the biomedical literature:
ecological host-parasite field observations are published primarily in
ecology journals (Web of Science, BioOne) that are not indexed by PubMed,
while clinical and biomedical pathogen papers dominate PubMed's open-access
content. The current performance ceiling at F1\,$\approx$\,0.87 is therefore
determined by data distribution rather than model capacity.

% ─────────────────────────────────────────────────────────────────────────────
\section{Discussion}
\label{sec:discussion}

\subsection{Soft Labels as the Primary Driver of Improvement}

The most actionable finding is that soft label supervision provides the
largest performance gain of any intervention tested: +0.052 F1, more than
the NER auxiliary task (+0.021) and more than switching from a single model
to an ensemble (+0.077 cumulatively, but requiring 3.3$\times$ inference cost).
The mechanism is consistent with the knowledge distillation literature
\citep{hinton2015distilling, tang2019distilling}: a teacher model that has
been exposed to diverse linguistic expressions of interactions assigns
calibrated probabilities to borderline sentences, encoding genuine semantic
uncertainty that hard labels cannot represent.

\subsection{NER as Zero-Cost Auxiliary Supervision}

The NER auxiliary task improves classification without requiring any
additional manually annotated data. BIO labels are derived automatically
from species gazetteers and interaction lexicons, making this a form of
distant supervision \citep{mintz2009distant}. The improvement mechanism is
indirect: NER pre-training updates the shared encoder to be more sensitive
to species names and interaction verbs, producing richer [CLS] token
representations for the classification head.

This approach is distinct from standard NER-augmented pipelines
\citep{habibi2017deep} in which entity detection and relation classification
are performed sequentially. Here, NER supervision modifies the encoder
representations used for classification; the NER output itself is not used
at inference time.

\subsection{Utility for BiotXplorer and Knowledge Graph Construction}

The classifier provides two concrete operational benefits for BiotXplorer.
First, as an upstream filter, it eliminates sentences that merely co-mention
species without describing an interaction, reducing index noise and improving
search relevance. Second, the calibrated probability output allows the
BiotXplorer platform to rank results by interaction confidence rather than
simple keyword match, enabling more nuanced retrieval.

For downstream knowledge graph construction, the multi-task model's NER
labels provide additional structure at no extra cost: the HOST/PATHOGEN
distinction produced by the NER head directly encodes the direction of the
knowledge graph edge (who interacts with whom), enabling the downstream
relation extractor to assign typed edges rather than undirected links.
This is particularly relevant for pathogen-host and parasite-host pairs
where edge direction is biologically meaningful.

\subsection{The PubMed Distribution Ceiling}

The distribution mismatch analysis identifies a concrete obstacle to
improvement: ecological host-parasite and predation interactions account
for 58\% of test-set positives but fewer than 15\% of training positives.
These sentence types cannot be reliably harvested from PubMed because the
relevant literature (field ecology, ecological entomology, parasitology
monographs) is predominantly published in journals not indexed by PubMed.

The path beyond F1\,$\approx$\,0.87 requires harvesting sentences from
non-PubMed ecological literature, re-running the Qwen3.5-122B annotation
pipeline, and updating the ensemble soft labels on the expanded corpus.
This is a data problem, not a modelling problem: the current architecture
can accommodate improved training data without modification.

\subsection{Limitations}

\begin{enumerate}
    \item \textbf{Test set size.} 100 sentences yields 95\% CIs of
    $\pm$0.06--0.07 F1. Most pairwise model comparisons are not individually
    significant; a test set of 500--1000 sentences is needed for
    publication-quality model selection.

    \item \textbf{Binary classification.} The classifier detects interactions
    but does not classify interaction type. Downstream relation extraction
    is needed to produce typed knowledge graph edges.

    \item \textbf{English only.} Training and evaluation are exclusively
    English-language.

    \item \textbf{Automated NER labels.} BIO labels are generated by gazetteers
    and lexicons rather than human annotators. The NER auxiliary loss trains
    the encoder against this automated signal, not a gold-standard annotation.
\end{enumerate}

% ─────────────────────────────────────────────────────────────────────────────
\section{Conclusion}

We have presented a multi-task BiomedBERT classifier for sentence-level
biotic interaction detection that achieves F1\,=\,0.868 on an ecological
test set. The model is built on three components: a two-stage teacher chain
that combines large-LLM binary annotation with ensemble soft probability
distillation; a multi-task NER auxiliary task that primes the shared encoder
at zero additional annotation cost; and a two-phase training protocol that
separates entity-level supervision from joint optimisation.

The central quantitative finding is that soft labels are the single most
impactful intervention in the pipeline, contributing +0.052 F1 on identical
architecture and data. The current performance ceiling at F1\,$\approx$\,0.87
is explained by a structural distribution mismatch between PubMed-sourced
training data and the ecological interaction types dominant in the test set,
pointing directly to the data collection strategy needed to break this ceiling.

The classifier is deployed in BiotXplorer
(\url{https://biotxplorer.sibils.org/}) as an upstream quality filter
and constitutes the first component of a knowledge graph construction
pipeline for biodiversity and host-pathogen interactions. Future work
will address: (1) harvesting from non-PubMed ecological literature to
expand training coverage of parasitism and host-interaction types;
(2) typed relation extraction to produce directed knowledge graph triples;
(3) integration with amplicon sequencing species detection for automated
interaction retrieval from biodiversity field surveys.

% ─────────────────────────────────────────────────────────────────────────────
\section*{Data and Code Availability}

Training code, evaluation scripts, and the test set are available
at the project repository. Model weights can be reproduced from the
provided training configuration.

\bibliographystyle{abbrvnat}
\bibliography{references}

\end{document}
